\chapter{Overview and Architecture}
\label{ch:overview}

\section{What OpenSG-2.0 is}
\label{sec:what-is}

OpenSG-2.0 is an implementation of Mechanics of Structure Genome (MSG)
homogenization written in JAX. You give it a \emph{structure gene} (SG) --- the
smallest piece of a structure that carries its heterogeneity --- and it returns
the macroscopic constitutive law of the structure that SG builds: a Timoshenko
beam $6\times6$, a plate ABD (or the shear-refined $8\times8$), or an equivalent
3-D solid $6\times6$. It also runs the same calculation backwards
(\emph{dehomogenization}), rebuilding the pointwise 3-D stress and strain inside
the SG from macroscopic beam or plate loads.

The repository is organised in four layers. Only the top layer knows where files
live on disk.

\begin{table}[htbp]
\centering
\caption{The four layers of OpenSG-2.0. Paths and \texttt{main} blocks exist
only in the top row.}
\label{tab:layers}
\small
\begin{tabular}{@{}llp{7.4cm}@{}}
\toprule
Layer & Package & Role \\
\midrule
user     & \texttt{examples/}        & YAML in $\rightarrow$ \texttt{.out} out; owns all paths \\
msg-shell& \texttt{src/opensg\_shell}& Reissner--Mindlin shell SGs: ring (cross-section), aperiodic segment, 3-D shell cell \\
msg-solid& \texttt{src/opensg\_solid}& solid SGs: plate ABD, beam KKT, 3-D solid law \\
FE core  & \texttt{src/fe\_jax}      & JAX FE architecture: basis/quadrature, periodic assembly map, jitted element kernels, element-by-element (EBE) Chebyshev-CG, sparse direct \\
\bottomrule
\end{tabular}
\end{table}

\subsection{The universal contract}
\label{sec:contract}

Every capability in both engines obeys the same contract: \textbf{a YAML file (or
a SwiftComp \texttt{.sc} file) goes in, and a timed \texttt{.out} file in the
SwiftComp \texttt{.K} layout comes back.} There is exactly one writer for that
file, \texttt{opensg\_solid.sg\_homo.write\_sc\_K}, and it is called from inside
the solver source, not from the example scripts, so no run can forget it. A
minimal session looks like this:

\begin{lstlisting}[style=opensg,caption={The whole contract: the .out file is written as a side effect of the call.},basicstyle=\ttfamily\footnotesize]
from opensg_shell import build_rm_bundle

B = build_rm_bundle("iea_s10_shell.yaml")   # writes iea_s10_shell_Timo.out
B["Timo"]                                    # the Timoshenko 6x6 as a numpy array
\end{lstlisting}

The file that comes back always has the same shape: an \texttt{OpenSG} banner
naming the model and the SG measure $\omega$ on the first line, the effective
stiffness matrix, its inverse (the compliance), an engineering-constants block
for 3-D laws only, and \texttt{Time taken} on the last line. Numbers are printed
in a $19$-character right-justified Fortran field
(\verb|%.7E| mantissa, signed three-digit exponent) so the file drops straight
into existing SwiftComp comparison tooling:

\begin{Verbatim}[frame=single,fontsize=\small,samepage=false]
 OpenSG msg-shell beam model [ext sh2 sh3 twist bend2 bend3]

 The Effective Timoshenko Beam Stiffness Matrix
 --------------------------------------------
     <6 numbers per row>
 ...
 The Effective Timoshenko Beam Compliance Matrix
 --------------------------------------------
 ...
 Time taken: 12.34 sec
\end{Verbatim}

The suffix on the output file names the model, so two different routes run on
the same input YAML do not silently overwrite each other
(Table~\ref{tab:outfiles}).

\begin{table}[htbp]
\centering
\caption{Output file names. There is one file per model; the \texttt{\_Timo} vs
\texttt{\_C3D} suffix disambiguates when two routes can read the same YAML.}
\label{tab:outfiles}
\small
\begin{tabular}{@{}llp{5.2cm}@{}}
\toprule
Model & Entry function & Output file \\
\midrule
msg-solid beam / plate / 3-D & \texttt{plate\_homo\_2d} & \texttt{<base>.out} \\
msg-shell beam (Timoshenko)  & \texttt{build\_rm\_bundle} & \texttt{<yaml>\_Timo.out}, \texttt{<yaml>\_ABDG.out} \\
msg-shell aperiodic segment  & \texttt{segment\_timo\_from\_3dyaml} & \texttt{<yaml>\_Timo.out} \\
msg-shell solid props, cross-section SG & \texttt{build\_solid\_bundle} & \texttt{<yaml>\_C3D.out}, \texttt{<yaml>\_ABDG.out} \\
msg-shell 3-D shell SG       & \texttt{shell\_sg3d} & \texttt{<yaml>\_C3D.out}, \texttt{<yaml>\_ABDG.out} \\
\bottomrule
\end{tabular}
\end{table}

\subsection{One command: \texttt{opensg}}
\label{sec:opensg-cli}

You do not have to know which engine owns a file in order to run it. The
package installs one dispatcher, \texttt{opensg}, alongside the two
engine-specific commands \texttt{opensg\_solid} and \texttt{opensg\_shell};
all three take the same one or two arguments:

\begin{lstlisting}[style=shellst]
opensg <sg.yaml>          # homogenization -- the default
opensg <sg.yaml> H        # homogenization, stated explicitly
opensg <sg.yaml> D        # dehomogenization: homogenize, then recover
\end{lstlisting}

\noindent
\texttt{opensg} (\texttt{src/opensg/cli.py}) is a pure dispatcher and prints
nothing of its own. It resolves the engine with
\texttt{opensg\_solid.sg\_mesh.resolve\_msg} --- the YAML's \texttt{msg:} key
(\texttt{shell} or \texttt{solid}) or, when that key is absent, the mesh
\emph{dialect} (\texttt{nodes}/\texttt{cells}/\texttt{mat\_id} is solid,
\texttt{nodes}/\texttt{elements}/\texttt{sets}/\texttt{sections}/\texttt{elementOrientations}
is shell) --- and hands the untouched argument list to
\texttt{opensg\_shell.cli.main} or \texttt{opensg\_solid.cli.main}. A
\texttt{msg:} that contradicts the dialect is an error, not a silent
override, and a file written before the key existed still runs. The
engine-specific commands remain and behave identically, as do
\texttt{python -m opensg\_solid <yaml>} and
\texttt{python -m opensg\_shell <yaml>}. Everything else --- the macro
model, the classical/shear-refined switch, the macro state for a recovery
--- lives in the YAML header, which is why the command line needs no flags.

Every run prints the same header, so the terminal states what was actually
read before any number appears:

\begin{Verbatim}[frame=single,fontsize=\small,samepage=false]
============================================================================
OpenSG -- a multiscale structural analysis tool based on the Mechanics of
Structure Genome (MSG), developed by the Multiscale Structural Mechanics
Group led by Prof. Wenbin Yu at Purdue University
============================================================================

 input     : /abs/path/to/iea_s10_shell.yaml
 msg       : shell
 SG dim    : 2D
 analysis  : homogenization
 macro model: beam, shear-refined

Timoshenko Beam Stiffness Matrix  [eps11 gam12 gam13 kappa1 kappa2 kappa3]:
[[ 2.76606e+10 ... ]]
Homogenization stored in iea_s10_shell_Timo.out
Time taken: 2.58 sec
\end{Verbatim}

Three lines of that header repay reading. \texttt{msg} is the engine that
was resolved, so a mis-tagged file is visible on the first line of output.
\texttt{SG dim} is the \textbf{ambient} dimension --- the space the SG
occupies, \texttt{node\_span\_dim}, not the intrinsic dimension of the
manifold --- so a Schwarz-P shell surface prints \texttt{3D} while a
cross-section ring and a 2-D solid cell both print \texttt{2D}.
\texttt{macro model} echoes \texttt{n\_model} and \texttt{refined}, and gains
a trailing \texttt{, aperiodic} when the header sets \texttt{aperiodic}. For
a dehomogenization the two closing lines are instead
\texttt{Local field files are computed and stored.} and a single
\texttt{Time taken}, covering the whole run.

\subsection{The cores hold no paths}
\label{sec:no-paths}

Neither \texttt{opensg\_shell} nor \texttt{opensg\_solid} contains a
\texttt{main} block, a hard-coded directory, or an assumption about where you
run from. Every path decision lives in the example script, which names its input
file and calls exactly one entry function on it. That is why the tutorials in
the later chapters can be reduced to two lines of user input and one call: to
run a different cross-section you change the string, not the library.

Two pieces of process-level setup do live in the package \texttt{\_\_init\_\_}
files, and both matter to you:

\begin{itemize}
\item \textbf{Double precision is load-bearing.} Both packages execute
      \texttt{jax.config.update("jax\_enable\_x64", True)} at import. Without
      it JAX silently degrades to float32 and the digits of the $6\times6$ are
      wrong. Import the package (\texttt{import opensg\_shell}) before you
      import anything else that touches JAX.
\item \textbf{A persistent compile cache.} \texttt{opensg\_solid/\_\_init\_\_.py}
      points JAX's compilation cache at \verb|$JAX_COMPILATION_CACHE_DIR|,
      defaulting to \verb|~/.cache/opensg_jax|, with
      \texttt{jax\_persistent\_cache\_min\_compile\_time\_secs} set to $0$. A
      cold run pays each XLA compile once; warm runs load the binaries back. The
      cache is wrapped in a \texttt{try}/\texttt{except} --- it is an
      optimisation, and the engine must run without it. The same file honours
      \verb|$FE_JAX_CORE|, which prepends an external \texttt{fe\_jax} checkout
      to \texttt{sys.path} to override the bundled copy at \texttt{src/fe\_jax}.
\end{itemize}

\section{The theory in one page}
\label{sec:theory-one-page}

Every pipeline in this manual is the same variational statement, specialised. It
is worth reading this section once, because the symbols used here are the
variable names used in the source: \texttt{Dee}, \texttt{Dhe}, \texttt{Dhh},
\texttt{V0}, \texttt{V1}, \texttt{omega}.

\subsection{The structure gene and the strain split}

A structure gene is the smallest piece of the structure that carries the
heterogeneity: a laminate stacking sequence through the thickness (a 1-D SG), a
cross-section contour or a 2-D unit cell (a 2-D SG), a lattice or TPMS unit cell
(a 3-D SG). The displacement field inside the SG is split into a macroscopic
part driven by the macro strains $\bar\varepsilon$ and a fluctuation (warping)
field $w$. In strain form,

\begin{equation}
\varepsilon \;=\; \Gamma_e\,\bar\varepsilon \;+\; \Gamma_h\,w ,
\label{eq:split}
\end{equation}

where $\Gamma_e$ is the operator that embeds the macro strains into the SG and
$\Gamma_h$ is the SG's own strain operator acting on the warping. The two are
not independent: $\Gamma_e$ is $\Gamma_h$ evaluated on the macro embedding, a
consistency requirement written down in
\texttt{Rules/gamma\_e\_gamma\_h\_consistency.md} and enforced by building both
from the same element kernel.

\subsection{The quadratic form}

Substituting~\eqref{eq:split} into the strain energy gives the quadratic form
the code assembles, in exactly the block names used throughout the source:

\begin{equation}
2U \;=\; \bar\varepsilon^{T}\mathbf{D}_{ee}\,\bar\varepsilon
      \;+\; 2\,w^{T}\mathbf{D}_{he}\,\bar\varepsilon
      \;+\; w^{T}\mathbf{D}_{hh}\,w .
\label{eq:energy}
\end{equation}

Read it as three integrals over the SG: $\mathbf{D}_{ee}$ is the macro-macro
block (what you would get with no warping at all, the rule-of-mixtures answer),
$\mathbf{D}_{hh}$ is the SG stiffness seen by the warping alone, and
$\mathbf{D}_{he}$ couples them.

\subsection{The zeroth-order solve}

Minimising~\eqref{eq:energy} over $w$, subject to the SG's boundary treatment
(periodicity, Dirichlet boundary data, or a rigid-body constraint that removes
the null space) gives one linear system with as many right-hand sides as there
are macro strains,

\begin{equation}
\mathbf{D}_{hh}\,\mathbf{V}_0 \;=\; -\,\mathbf{D}_{he},
\label{eq:v0}
\end{equation}

and the effective law follows by substituting the minimiser back:

\begin{equation}
\mathbf{D}_{\mathrm{eff}} \;=\; \mathbf{D}_{ee} + \mathbf{V}_0^{T}\mathbf{D}_{he},
\qquad
\mathbf{C} \;=\; \mathbf{D}_{\mathrm{eff}}/\omega .
\label{eq:deff}
\end{equation}

Here $\omega$ is the SG measure that stays in the model --- the thickness of a
1-D laminate SG, the area of a cross-section SG, the volume of a 3-D solid cell,
the perimeter (or, in three dimensions, the midsurface area) of a shell SG. The
banner of every \texttt{.out} states $\omega$ because a bare matrix is
ambiguous: a whole-matrix ratio that is uniform across all entries (exactly
$0.5000$, say) is a normalisation difference, not a physics discrepancy.

$\mathbf{V}_0$ is stored, not discarded. Its columns are the warping fields per
unit macro strain, and they are what the dehomogenization reads back.

Equation~\eqref{eq:deff} is the \textbf{zeroth-order} answer: the plate ABD, the
Euler--Bernoulli $4\times4$, the equivalent 3-D solid law. It contains no
transverse shear, because the classical macro models it corresponds to have
none.

\subsection{The first-order step: transverse shear}

Recovering transverse shear --- the Reissner--Mindlin $G$ block of a plate, the
$GA_{12}$ and $GA_{13}$ entries of a Timoshenko $6\times6$ --- takes a
\textbf{first-order} step: a second solve for $\mathbf{V}_1$ against a
right-hand side built from $\mathbf{V}_0$, followed by an energy transformation.
Two functions do this, and they are shared by every beam route in the
repository:

\begin{itemize}
\item \texttt{prepare\_v1\_rhs(V0, Dhl, Dll, Dle\_dense, Psi, Dc)} forms the
      first-order right-hand side and the auxiliary blocks
      $\mathbf{V}_0^{T}\mathbf{D}_{ll}\mathbf{V}_0$,
      $\mathbf{D}_{hl}\mathbf{V}_0$ and
      $\mathbf{D}_{hl}^{T}\mathbf{V}_0 + \mathbf{D}_{le}$, projecting the
      right-hand side onto the range of $\mathbf{D}_{hh}$ so the singular
      system stays consistent.
\item \texttt{finalize\_v1\_and\_compute\_deff(V1s\_raw, V0, D\_eff, V0DllV0,
      DhlV0, DhlTV0Dle, Psi, Dc)} projects $\mathbf{V}_1$, symmetrises
      $\mathbf{D}_{\mathrm{eff}}$, and performs the generalized-Timoshenko
      transformation of the energy blocks $(A, B, C, D)$ into the sorted
      $6\times6$.
\end{itemize}

Both live in \texttt{opensg\_shell/fe\_jax/msg\_solver.py} and are
\texttt{@jax.jit}-decorated. The output ordering of the $6\times6$ is
\texttt{[EA, GA12, GA13, GJ, EI2, EI3]}, i.e.\ the strain measures
$[\varepsilon_{11}\;\;\gamma_{12}\;\;\gamma_{13}\;\;\kappa_1\;\;\kappa_2\;\;\kappa_3]$,
which is the VABS diagonal order.

The symmetrisation inside \texttt{finalize\_v1\_and\_compute\_deff} is not
cosmetic. $\mathbf{D}_{\mathrm{eff}}$ is a Hessian and therefore symmetric by
construction, but a Dirichlet-constrained tapered-segment solve produces a small
($2$ to $6$ percent) antisymmetric artifact, because with $\mathbf{V}_0$ pinned
at the end rings $\mathbf{D}_{ee}+\mathbf{V}_0^{T}\mathbf{D}_{he}$ is no longer
the exact symmetric Schur complement. Left alone, that artifact propagates
through the inverse into the antisymmetric shear--bending coupling ($C_{26}$,
$C_{35}$) and inflates it by up to $78$ percent on anisotropic tapers while
leaving the symmetric terms untouched. Enforcing the invariant is a no-op for a
prismatic SG.

\subsection{One factorization for two solves}

Equations~\eqref{eq:v0} and the $\mathbf{V}_1$ system share the same operator.
Every beam route therefore factorizes once and back-substitutes twice. In
\texttt{sg\_homo.ring\_indep} the augmented KKT matrix
\[
\begin{bmatrix}
\mathbf{D}_{hh} & \mathbf{D}_c\\[2pt]
\mathbf{D}_c^{T} & \mathbf{0}
\end{bmatrix}
\begin{Bmatrix}\mathbf{V}\\ \lambda\end{Bmatrix}
=
\begin{Bmatrix}\mathbf{R}\\ \mathbf{0}\end{Bmatrix}
\]
is built once, passed to \texttt{scipy.linalg.lu\_factor}, and then
\texttt{lu\_solve} is called with $\mathbf{R}=-\mathbf{D}_{he}$ for
$\mathbf{V}_0$ and with the \texttt{prepare\_v1\_rhs} output for
$\mathbf{V}_1$. In the tapered-segment route the same idea appears as
\texttt{sg\_assembly.dirichlet\_factor} followed by two
\texttt{dirichlet\_solve\_fac} calls; the second call falls back to a fresh
factorization only if the $\mathbf{V}_1$ Dirichlet set differs from the
$\mathbf{V}_0$ one. In the msg-solid beam engine the equivalent is
\texttt{solve\_fluctuation\_field}, which returns the assembled KKT matrix
alongside $\mathbf{V}_0$ precisely so the caller can reuse it; there the
constraint rows are scaled by $10^{8}$ before assembly to keep the saddle-point
system well conditioned, and empty rows are given a unit diagonal with a zeroed
right-hand side.

\subsection{Dehomogenization: running the chain backwards}

Given macro strains (or macro forces, converted through
$\bar\varepsilon = \mathbf{C}^{-1}\mathbf{F}$), the warping is rebuilt from the
stored $\mathbf{V}_0$ and $\mathbf{V}_1$, equation~\eqref{eq:split} is evaluated
pointwise, and the material law of each element turns that strain into 3-D
stress. Nothing is re-solved: homogenize once, recover as often as you like. For
the shell engine this is a \emph{two-step} recovery --- step 1 rebuilds the six
shell strains plus the two wall transverse shears
$[\varepsilon_{11}\;\varepsilon_{22}\;2\varepsilon_{12}\;\kappa_{11}\;\kappa_{22}\;2\kappa_{12}]$
and $[2\gamma_{13}\;2\gamma_{23}]$ on the ring, and step 2 pushes them through
the through-thickness plate SG to get ply-level 3-D stress. The entry points are
\texttt{opensg\_shell.stress\_at\_points} and
\texttt{opensg\_shell.disp\_at\_points} for the shell engine, and
\texttt{opensg\_solid.sg\_dehom.dehom\_fields} / \texttt{plate\_dehom\_2d} for
the solid engine.

\section{Module map}
\label{sec:module-map}

Both packages are split file-per-concern with the same naming, so a routine you
know in one engine is in the same-named file in the other.

\begin{table}[htbp]
\centering
\caption{\texttt{src/opensg\_shell} --- the MSG shell engine.}
\label{tab:mod-shell}
\small
\begin{tabular}{@{}lp{9.1cm}@{}}
\toprule
File & Concern \\
\midrule
\texttt{sg\_mesh.py} &
SG input loaders and mesh utilities: 1-D ring YAML loaders
(\texttt{load\_ring}, \texttt{load\_ring\_ref}) with explicit laminate-reference
conventions; 3-D segment loading and topological boundary extraction
(\texttt{extract}) into solver \texttt{.npz} bundles; orientation-frame numeric
reports and PNG prechecks. \\
\texttt{sg\_materials.py} &
Section laws: the per-section wall plate ABD $6\times6$ plus the $2\times2$
transverse-shear $G$ (\texttt{\_material\_by\_section}), and the per-station
wall-law emitter \texttt{emit\_station\_abd} / \texttt{load\_station\_abd} that
writes one windIO-style YAML of $8\times8$ Reissner--Mindlin plate laws per
cross-section. \\
\texttt{sg\_assembly.py} &
Shell element operators and assembly: the production 6-DOF drilling element, the
5-DOF general taper element, the prismatic MITC4 element, MITC assumed-strain
tying, the \texttt{compute\_k22} / \texttt{compute\_kg} initial-curvature
corrections, the Dirichlet factorizations
(\texttt{dirichlet\_factor}, \texttt{dirichlet\_solve\_fac}), and the six-block
MSG energy assembly (\texttt{assemble\_segment\_indep},
\texttt{assemble\_constraint}). \\
\texttt{sg\_periodicity.py} &
Periodic local-to-global sparse assembly map for a shell mesh, $6$ DOF per node
$(w_1,w_2,w_3,\omega_1,\omega_2,\omega_3)$ instead of the solid engine's three
translations. The tie is a plain index map; no transformation matrix is formed. \\
\texttt{sg\_homo.py} &
Homogenization drivers: \texttt{ring\_indep} (ring $\rightarrow$ Timoshenko
$6\times6$), \texttt{build\_rm\_bundle} (the packaged bundle the dehom
consumes), \texttt{ring\_solid} / \texttt{build\_solid\_bundle}
(cross-section $\rightarrow$ equivalent 3-D solid), \texttt{shell\_sg3d} (3-D
shell SG $\rightarrow$ \texttt{C3D}), \texttt{segment\_timo\_from\_3dyaml}
(aperiodic / tapered segment), \texttt{elastic\_constants}. \\
\texttt{sg\_dehom.py} &
Reissner--Mindlin two-step dehomogenization and recovery:
\texttt{stress\_at\_points}, \texttt{disp\_at\_points}, built on the same
element operators as the homogenization so it is the energy-consistent adjoint. \\
\texttt{sg\_junction.py} &
Level-2 junction law: local corner and micro-cell solves
(\texttt{corner\_micro\_law} for L/T/X topologies, \texttt{microcell\_law} for a
periodic mini-lattice) that correct the shell lattice for sub-shell junction
physics. \\
\texttt{fe\_jax/} &
The OpenSG-authored \texttt{msg\_*} kernels: \texttt{msg\_materials}
(ABD and $6\times6$ tables, reference shifts), \texttt{msg\_transverse\_shear}
(wall $G$, the $8\times8$ plate law), \texttt{msg\_mesh},
\texttt{msg\_rm}/\texttt{msg\_rm\_timo} (RM operators and the 5-DOF Timoshenko
assembly), \texttt{msg\_solver} (KKT solve, $\mathbf{V}_1$ RHS, Timoshenko
finalization), \texttt{msg\_hermite} (the Hermite $C^1$ Kirchhoff--Love
thin-walled pipeline), \texttt{msg\_dehom}, \texttt{orient\_plot}. \\
\texttt{helper/} &
Format converters (\texttt{ff\_to\_yaml.py}). \\
\bottomrule
\end{tabular}
\end{table}

\begin{table}[htbp]
\centering
\caption{\texttt{src/opensg\_solid} --- the general structure-gene engine.}
\label{tab:mod-solid}
\small
\begin{tabular}{@{}lp{9.1cm}@{}}
\toprule
File & Concern \\
\midrule
\texttt{sg\_materials.py} &
The $6\times6$ stiffness tables, per material and per element:
\texttt{build\_material\_C} and \texttt{get\_heterogeneous\_C\_matrix}, plus the
in-plane rotation. Voigt order everywhere is SwiftComp
$(xx, yy, zz, yz, xz, xy)$. \\
\texttt{sg\_mesh.py} &
SG input parsing and mesh utilities: \texttt{load\_sg\_input} (a SwiftComp
\texttt{.sc} or the helper's YAML), \texttt{\_cell\_basis} (SG dimension plus
nodes-per-element $\rightarrow$ basix cell type and basis degree), and
\texttt{plot\_sg\_mesh}, which draws the actual mesh coloured by material. \\
\texttt{sg\_assembly.py} &
Element kernels, periodic assembly, preconditioners and linear solvers:
\texttt{calculate\_RHS\_and\_Ke\_batch\_periodic},
\texttt{compress\_periodic\_cells\_jax}, \texttt{compute\_block\_inv\_diag},
\texttt{apply\_block\_precond}, \texttt{apply\_chebyshev\_precond},
\texttt{ebe\_jacobian\_product\_periodic},
\texttt{estimate\_max\_eigenvalue}, \texttt{\_sparse\_direct\_solve},
\texttt{solve\_fluctuation\_field},
\texttt{plate\_ladder\_element\_blocks}. \\
\texttt{sg\_homo.py} &
Homogenization drivers and the one public entry \texttt{plate\_homo\_2d}: the
periodic CG pipeline (\texttt{full\_homogenization\_pipeline},
\texttt{\_homo\_direct}) for \texttt{n\_model} 2 and 3, the beam KKT path
(\texttt{\_beam\_homo\_kkt}) for \texttt{n\_model} 1, the Reissner--Mindlin
first-order plate ladder (\texttt{plate\_shear\_ladder}), and the single output
writer \texttt{write\_sc\_K}. \\
\texttt{sg\_dehom.py} &
Dehomogenization: the Gauss-point recovery kernel for \texttt{n\_model} 2 and 3,
the generalized-Timoshenko beam recovery chain for \texttt{n\_model} 1, the
public \texttt{dehom\_fields} and \texttt{plate\_dehom\_2d}, and the Gauss
exports (\texttt{export\_gauss}). The batched recovery is wrapped in a
module-level \texttt{jit} so repeated calls in one process do not re-trace. \\
\texttt{rm\_plate\_1D/} &
The 1-D through-thickness plate stack used by both engines:
\texttt{segment\_plate.plate\_sg\_yaml} (build the through-thickness SG from a
stacking sequence), \texttt{msg\_rm\_plate.rm\_plate\_msg} (the $8\times8$ ABDG
and the warping ladders), \texttt{rm\_dehom} (ply-level recovery),
\texttt{rm\_homo}, \texttt{cli\_rm\_plate}. \\
\bottomrule
\end{tabular}
\end{table}

The shared FE core \texttt{src/fe\_jax} is not an OpenSG concern but a general
JAX finite-element architecture. The pieces the two engines actually import are
\texttt{basis\_quadrature} (\texttt{FiniteElementType}, \texttt{get\_quadrature},
\texttt{eval\_basis\_and\_derivatives}, all backed by basix),
\texttt{setup.mesh\_to\_jax} and
\texttt{setup.mesh\_to\_periodic\_sparse\_assembly\_map},
\texttt{solve\_cg} (a conjugate-gradient solver that returns convergence
information), \texttt{linear\_elasticity}, and \texttt{profiling}.

\section{The pipelines}
\label{sec:pipelines}

Four routes are described here in the order a new user meets them. Each is a
sequence of calls that the entry function makes on your behalf; you only ever
call the entry function. Figure~\ref{fig:pipelines} puts all four side by side:
one input file, one entry function, one \texttt{.out} file per branch.

\begin{figure}[htbp]
\captree{
  \node[funcnode] (f1) {\texttt{sg\_homo.build\_rm\_bundle}\\ \texttt{sg\_mesh.load\_ring\_ref} (it calls\\ \texttt{\_material\_by\_section} and\\ \texttt{compute\_k22}) $\rightarrow$ \texttt{rm\_plate\_msg}\\ $\rightarrow$ \texttt{sg\_homo.ring\_indep}};
  \node[inputnode, left=5mm of f1] (r1) {1-D shell SG ring YAML\\ \texttt{iea\_s10\_shell.yaml}};
  \node[outnode, right=5mm of f1] (o1) {\texttt{iea\_s10\_shell\_Timo.out}\\ \texttt{iea\_s10\_shell\_ABDG.out}\\ the Timoshenko $6\times6$};

  \node[funcnode, below=8mm of f1] (f2) {\texttt{sg\_homo.segment\_timo\_from\_3dyaml}\\ \texttt{sg\_mesh.extract} end rings $\rightarrow$\\ \texttt{ring\_indep} per side $\rightarrow$\\ \texttt{\_boundary\_dirichlet} $\rightarrow$\\ \texttt{dirichlet\_factor}};
  \node[inputnode, left=5mm of f2] (r2) {3-D shell segment YAML\\ \texttt{BAR\_URC\_numEl\_52\_}\\ \texttt{segment\_12.yaml}};
  \node[outnode, right=5mm of f2] (o2) {\texttt{BAR\_URC\_numEl\_52\_}\\ \texttt{segment\_12\_Timo.out}\\ per unit length, $L$ in the banner};

  \node[funcnode, below=8mm of f2] (f3) {\texttt{sg\_homo.shell\_sg3d}\\ one wall law for every element,\\ junction-edge count, periodic or\\ aperiodic boundary, one \texttt{splu}};
  \node[inputnode, left=5mm of f3] (r3) {3-D shell SG YAML\\ \texttt{schwarz\_p\_3Dshell.yaml}};
  \node[outnode, right=5mm of f3] (o3) {\texttt{schwarz\_p\_3Dshell\_C3D.out}\\ \texttt{schwarz\_p\_3Dshell\_ABDG.out}\\ the equivalent solid $6\times6$,\\ per unit-cell volume};

  \node[funcnode, below=8mm of f3] (f4) {\texttt{sg\_homo.plate\_homo\_2d}\\ \texttt{load\_sg\_input} $\rightarrow$ \texttt{\_cell\_basis} $\rightarrow$\\ \texttt{get\_heterogeneous\_C\_matrix} $\rightarrow$\\ \texttt{\_beam\_homo\_kkt} (\texttt{n\_model} 1) or\\ \texttt{\_homo\_direct} (\texttt{n\_model} 2, 3)};
  \node[inputnode, left=5mm of f4] (r4) {solid SG, a SwiftComp \texttt{.sc}\\ or a YAML:\\ \texttt{RHC\_SW\_2UC\_45.yaml}};
  \node[outnode, right=5mm of f4] (o4) {\texttt{RHC\_SW\_2UC\_45.out}\\ \texttt{n\_model} 1 named Euler-\\ Bernoulli Beam or Timoshenko\\ Beam; 2 Classical Plate or\\ Reissner-Mindlin; 3 Cauchy\\ Continuum};

  \draw[capedge] (r1) -- (f1); \draw[capedge] (f1) -- (o1);
  \draw[capedge] (r2) -- (f2); \draw[capedge] (f2) -- (o2);
  \draw[capedge] (r3) -- (f3); \draw[capedge] (f3) -- (o3);
  \draw[capedge] (r4) -- (f4); \draw[capedge] (f4) -- (o4);
}
\caption{The four pipelines, one row each: the three msg-shell routes and the
         msg-solid driver whose \texttt{n\_model} argument selects the macro
         model. Yellow is the file you edit, blue the entry function with the
         calls it makes in order, green the \texttt{.out} file it leaves
         behind. All four \texttt{.out} files are written by the one writer
         \texttt{opensg\_solid.sg\_homo.write\_sc\_K}, which is why they share
         a single layout; the \texttt{\_ABDG.out} companion --- written by
         every route that reduces a layup to a wall plate law, the
         cross-section rings and the 3-D shell SG alike --- is the one
         exception, written by \texttt{write\_abdg\_out}.}
\label{fig:pipelines}
\end{figure}

\subsection{msg-shell: 1-D ring $\rightarrow$ Timoshenko $6\times6$}
\label{sec:pipe-ring}

This is the composite blade cross-section route. The input is a 1-D shell SG
YAML: nodes and two-node elements tracing the wall contour,
\texttt{elementOrientations}, \texttt{sections} carrying a layup per element set,
and \texttt{materials}. The entry function is
\texttt{build\_rm\_bundle(shell\_yaml, ref=None, shear="mitc4\_g23",
g\_source=None)}, and it proceeds as follows.

\begin{enumerate}
\item \texttt{sg\_mesh.load\_ring\_ref} reads the YAML and returns the contour
      arrays --- nodal coordinates \texttt{rx}, connectivity \texttt{cells},
      the per-element section index \texttt{rsub}, the element normals
      \texttt{re3}, the axial index \texttt{ax} and the two cross-section
      indices \texttt{cross}. The laminate reference surface comes from the
      YAML's own \texttt{reference} field (absent $\rightarrow$ \texttt{center},
      the mid-surface), so homogenization and recovery cannot disagree; pass
      \texttt{ref} explicitly only to override. The reference is converted once
      into a thickness fraction \texttt{frac}
      (\texttt{center} $\rightarrow 0.5$, \texttt{oml} $\rightarrow 0.0$,
      \texttt{oml\_flip} and \texttt{iml} $\rightarrow 1.0$) and that single
      number then drives the ring laminate reference, the plate-SG $z$
      reference, the emitted ABD and the recovery depth conversion.
\item \texttt{sg\_materials.\_material\_by\_section} builds, per section, the
      wall plate ABD $6\times6$ at the chosen reference and the $2\times2$
      transverse-shear $G$.
\item The wall $G$ of every section is then replaced by
      \texttt{rm\_plate\_msg}, the MSG/VAM second-order-energy (Yu-2002
      least-squares) projection in \texttt{opensg\_solid.rm\_plate\_1D}, run
      at the same \texttt{fraction=frac}. \textbf{This is the only route and
      there is no switch.} The retired \texttt{g\_source} selector is still
      accepted by the signature so existing callers keep working, but any
      value other than \texttt{msg} raises a \texttt{DeprecationWarning} and
      is \emph{ignored} (\texttt{sg\_materials.check\_g\_source}), and an
      emitted ABD cache records \texttt{g\_source: msg}. The
      complementary-energy (Whitney-1973) shear flow survives only as the
      internal numerical fallback on the rare laminate whose MSG $U^{*}$ fit
      is not positive definite --- a guard, not a user-selectable
      alternative. The per-section laws are written to
      \texttt{<yaml>\_ABDG.out}.
\item \texttt{sg\_assembly.compute\_k22} computes the hoop curvature per edge,
      the initial-curvature correction that makes a curved wall stiffer than a
      faceted approximation of it.
\item \texttt{sg\_homo.ring\_indep} assembles the ring as a one-element-deep
      strip of quads: the 6-DOF drilling element
      $(w_1,w_2,w_3,\theta_1,\theta_2,\omega_3)$ with an element-wise Lagrange
      multiplier enforcing the drilling constraint, and MITC tying of
      $\gamma_{23}$ only. (Under span invariance $\gamma_{13}$ is algebraic in
      the directors and carries no fluctuation gradient, so only $\gamma_{23}$
      pairs a differentiated displacement with a rotation.) The six energy
      blocks \texttt{Dhh, Dhe, Dee, Dhl, Dll, Dle} come out of
      \texttt{assemble\_segment\_indep} and the drilling rows
      \texttt{Gc, Gl, Ge} out of \texttt{assemble\_constraint}; all are divided
      by the artificial strip depth $h$.
\item The KKT matrix is factorized once; $\mathbf{V}_0$ and $\mathbf{V}_1$ are
      obtained by two back-substitutions, and
      \texttt{finalize\_v1\_and\_compute\_deff} returns the generalized-Timoshenko
      $6\times6$, symmetrised.
\item \texttt{write\_sc\_K} writes \texttt{<yaml>\_Timo.out}, and the returned
      bundle carries \texttt{Timo}, \texttt{V0}, \texttt{V1}, the contour
      geometry, the strip, the per-element layup names and the layup and
      material databases --- everything the two-step dehomogenization needs.
\end{enumerate}

The shipped example is
\texttt{examples/OpenSG\_shell/1\_get\_beam\_props\_from\_shell\_cross\_section},
whose input is \texttt{iea\_s10\_shell.yaml} (the IEA-22 blade at $r/R=0.2$):

\begin{lstlisting}[style=shellst]
cd examples/OpenSG_shell/1_get_beam_props_from_shell_cross_section
opensg iea_s10_shell.yaml
\end{lstlisting}

\noindent
The YAML carries \texttt{n\_model: 1}, \texttt{refined: 1} and
\texttt{msg: shell}, so that one command runs the Timoshenko route and writes
\texttt{iea\_s10\_shell\_Timo.out} (the Timoshenko stiffness and compliance)
and \texttt{iea\_s10\_shell\_ABDG.out} (the per-section wall laws). The
printed diagonal is $[EA, GA_2, GA_3, GJ, EI_2, EI_3]$. The shipped driver
\texttt{beam\_homo\_shell.py} makes the same call through
\texttt{build\_rm\_bundle} and additionally draws the two figures the CLI
does not --- \texttt{iea\_s10\_mesh.png} (the contour coloured by layup) and
\texttt{iea\_s10\_orient.png} (the compulsory $e_1/e_2/e_3$ orientation
figure, emitted by \texttt{auto\_emit}) --- so run the script when you want
those:

\begin{lstlisting}[style=shellst]
python beam_homo_shell.py
\end{lstlisting}

\subsection{msg-shell: 3-D segment $\rightarrow$ aperiodic Timoshenko $6\times6$}
\label{sec:pipe-segment}

The ring pipeline assumes the beam is prismatic: periodic along the axis, so one
cross-section suffices. A tapered segment is not periodic, so the periodicity
condition has to be replaced. The entry function is
\texttt{segment\_timo\_from\_3dyaml(seg\_yaml, workdir=None,
lam\_space="elem", shear="full")} and the input is a 3-D shell segment YAML
(nodes, quads, sections, materials).

\begin{enumerate}
\item \texttt{sg\_mesh.extract} finds the two end cross-sections
      \emph{topologically}: a mesh edge used by exactly one quad is a free edge,
      and the connected components of the free-edge graph are the end rings.
      Each is written out as a standalone 1-D SG YAML (so you can inspect or
      re-run it on its own) together with an orientation PNG, and the whole
      bundle --- segment coordinates, cells, subdomains, the three element frame
      vectors, and the node-to-segment map --- is saved to
      \texttt{<base>\_boundaries.npz}.
\item Each end ring is solved as a ring in its own right by
      \texttt{ring\_indep}, giving \texttt{C6}, \texttt{V0} and \texttt{V1} per
      side.
\item The full 3-D quad mesh is assembled with the same
      \texttt{assemble\_segment\_indep} and \texttt{assemble\_constraint}, with
      the element hoop curvature \texttt{k22\_e} and geometric curvature
      \texttt{kg\_e} computed from the element centroids and frames. The
      drilling multipliers are appended as extra rows and columns by
      \texttt{\_augment\_drilling}, giving a saddle-point system with zero
      penalty.
\item \texttt{\_boundary\_dirichlet} maps the ring warping fields onto the
      segment's end nodes. This is the boundary condition that replaces
      periodicity: $\mathbf{V}_0 = \mathbf{V}_{0,\mathrm{ring}}$ and
      $\mathbf{V}_1 = \mathbf{V}_{1,\mathrm{ring}}$ on both end sections.
\item \texttt{dirichlet\_factor} factorizes the constrained interior once;
      $\mathbf{V}_0$ and then $\mathbf{V}_1$ are solved against it. The
      Euler--Bernoulli block is
      $A_{EB} = (\mathbf{D}_{ee} + \mathbf{V}_0^{T}\mathbf{D}_{he})/L$ with $L$
      the axial extent of the nodes, and every energy block handed to
      \texttt{finalize\_v1\_and\_compute\_deff} is likewise divided by $L$, so
      the result is a per-unit-length stiffness.
\item \texttt{<yaml>\_Timo.out} is written with the segment length in its
      banner; the returned dict carries \texttt{S6}, \texttt{C6L}, \texttt{C6R},
      \texttt{L} and \texttt{solve\_time}.
\end{enumerate}

Two acceptance checks come with this route. On a \emph{prismatic} segment the
segment $6\times6$ must reproduce its own boundary ring $6\times6$; measured on
an isotropic and a $\pm45$ tube this identity holds to
$\le 5\times10^{-5}$ relative. On a genuinely tapered segment no identity holds,
and the check is instead that each diagonal term of the segment lies between the
two ring values --- the ratio $(S_{ii}-R_{ii})/(L_{ii}-R_{ii})$ printed by the
example must fall in $[0,1]$. The shipped case is a BAR-URC 100\,m blade
segment:

\begin{lstlisting}[style=shellst]
cd examples/OpenSG_shell/5_taper_segment_aperiodic_boundaries
python run_bar_urc_segment.py
\end{lstlisting}

\noindent
which reads \texttt{BAR\_URC\_numEl\_52\_segment\_12.yaml} and writes
\texttt{BAR\_URC\_numEl\_52\_segment\_12\_Timo.out},
\texttt{BAR\_URC\_numEl\_52\_segment\_12\_boundaries.npz}, the two boundary
YAMLs \texttt{boundary\_..\_L.yaml} and \texttt{boundary\_..\_R.yaml}, three
orientation PNGs, and the comparison table \texttt{bar\_urc\_segment12.dat}.

\subsection{msg-shell: 3-D shell SG $\rightarrow$ equivalent solid $C_{3D}$}
\label{sec:pipe-shell3d}

Here the SG is a surface in three dimensions --- a triply periodic minimal
surface, a lattice of sheets --- and the answer wanted is the $6\times6$ of the
equivalent homogeneous 3-D solid. The entry function is
\texttt{shell\_sg3d(yaml\_path, omega=None, drill\_pen=1.0e-3,
g\_source=None, solver="direct", boundary=None, shear="mitc")}.
\texttt{boundary=None} means \emph{read it from the YAML header}, where the
optional \texttt{aperiodic: 1} key is the opt-out --- the same key
\texttt{opensg\_solid.plate\_homo\_2d} reads, so the two CLIs share one
header contract; an explicit argument always wins.

\begin{enumerate}
\item The YAML's nodes, elements and \texttt{elementOrientations} are read; the
      element normal $e_3$ is columns 7 through 9 of the orientation rows.
      Triangles and quads may be mixed: they are assembled as two separate
      element batches (MITC3 and MITC4) scattered into the same global
      system, never padded into one another.
\item The single wall law (ABD at the mid-surface reference plus the $2\times2$
      $G$, again from \texttt{rm\_plate\_msg}, the only $G$ route) is built
      once and applied to every element.
\item Junction lines are counted: any shell edge shared by more than two
      elements. The count appears in the \texttt{.out} banner, because it tells
      you how much of the cell's stiffness comes from junction physics the
      plain shell lattice does not resolve.
\item The boundary treatment. With \texttt{boundary="periodic"}, the default, the
      periodic map ties opposite faces, edges and corners through the sparse
      assembly map, and the three rigid translations are removed by three
      area-weighted Lagrange rows. With \texttt{boundary="aperiodic"} the macro
      (affine) field is mapped onto the boundary instead, prescribing zero
      \emph{translational} fluctuation $w_1=w_2=w_3=0$ on every node of the
      bounding-box faces while the rotational DOFs stay free. Clamping the
      rotations as well over-stiffens the cut edges; freeing them was measured
      to move $C_{11}$ from $+12$ to $+7$ percent. The aperiodic result is a
      kinematic upper bound on a single cell and converges to the periodic one
      as the SG grows.
\item The batched $\Gamma_h$ and $\Gamma_e$ operators are evaluated at the Gauss
      points, contracted into the sparse $\mathbf{K}$, $\mathbf{D}_{he}$ and
      $\mathbf{D}_{ee}$, with a drilling stabilisation of
      \texttt{drill\_pen} $\times\,A_{11}$ on the $\omega_3$ DOF. One
      \texttt{scipy.sparse.linalg.splu} factorization then serves all six macro
      strain columns.
\item $\mathbf{D}_{\mathrm{eff}} = \mathbf{D}_{ee} +
      \mathbf{V}_0^{T}\mathbf{D}_{he}$ is symmetrised and divided by $\omega$.
      The default $\omega$ is the \emph{volume the equivalent continuum
      occupies}, i.e.\ the node bounding-box volume measured from the mesh
      --- the same convention \texttt{opensg\_solid} uses, and the same cell
      the periodic assembly map ties --- so the returned \texttt{C3D}, the
      printed matrix and \texttt{<yaml>\_C3D.out} are one and the same
      matrix, directly comparable with a solid \texttt{.K}. Pass
      \texttt{omega} explicitly (or set the YAML \texttt{omega:} header key)
      only when the continuum occupies something else; the midsurface
      \emph{surface area}, for a per-unit-area wall law, is returned
      alongside as \texttt{surface\_area}.
\item The step-1 record --- the $8\times8$ Reissner--Mindlin wall plate law
      and its compliance, one block per section --- is written to
      \texttt{<yaml>\_ABDG.out} by the same \texttt{write\_abdg\_out} the
      cross-section routes use.
\end{enumerate}

Because the wall law is a Reissner--Mindlin plate law, the sheet thickness is a
parameter of the material model rather than of the mesh: the same surface mesh
serves any thickness, which is what makes a thickness sweep cheap.
Figure~\ref{fig:schwarz} shows the Schwarz-P unit cell that the shipped example
uses.

\begin{figure}[htbp]
\centering
\includegraphics[width=0.55\textwidth]{\FIGDIR/schwarz_p_mesh.png}
\caption{The Schwarz-P triply periodic minimal surface meshed as a 3-D shell
structure gene. The surface carries a Reissner--Mindlin wall law, so sheet
thickness is a material parameter and one mesh serves an entire thickness
sweep. Produced by \texttt{schwarz\_solid\_props.py} in
\texttt{examples/OpenSG\_shell/4\_get\_solid\_props\_from\_shell\_3D\_SG}.}
\label{fig:schwarz}
\end{figure}

\begin{lstlisting}[style=shellst]
cd examples/OpenSG_shell/4_get_solid_props_from_shell_3D_SG
python make_schwarz_yaml.py          # generates schwarz_p_3Dshell.yaml
opensg schwarz_p_3Dshell.yaml        # homogenizes it
\end{lstlisting}

\noindent
The first command is the mesh generator and has no CLI equivalent; the second
is the plain homogenization, which the YAML's own header
(\texttt{n\_model: 3}, \texttt{msg: shell}) fully describes. The run prints
the junction-edge count, the number of DOF and the solve time, then the
$6\times6$; it writes \texttt{schwarz\_p\_3Dshell\_C3D.out} and
\texttt{schwarz\_p\_3Dshell\_ABDG.out}. The driver
\texttt{schwarz\_solid\_props.py} makes the same call through
\texttt{shell\_sg3d} and additionally writes \texttt{schwarz\_p\_mesh.png}.
The companion script
\texttt{compare\_boundary\_modes.py} runs the same cell under both boundary
treatments and tabulates them in \texttt{boundary\_compare.dat}. This route is
validated against its own solid counterpart: for the Schwarz-P cell the shell SG
sits within $-3$ percent on the normal moduli, $-6$ percent on the shears and
$+0.3$ percent on Poisson's ratio relative to the SwiftComp-exact msg-solid
answer.

\subsection{msg-solid: one driver, three macro models}
\label{sec:pipe-solid}

The msg-solid engine has a single public entry,

\begin{lstlisting}[style=opensg,caption={The msg-solid driver signature: n\_model selects which macro strain set is built.},basicstyle=\ttfamily\footnotesize]
from opensg_solid.sg_homo import plate_homo_2d

r = plate_homo_2d(sc_path,               # a .sc or .yaml SG file
                  material_param=None,   # (n_mat, 9) engineering override
                  angles=None,           # deg per material
                  n_model=None,          # 1 beam, 2 plate, 3 3-D solid
                  refined=None,          # 0 classical, 1 shear-refined
                  workdir=None,
                  elem_rotation=None,    # (E, 9) per-element direction cosines
                  solver="direct",       # "direct" | "cg"
                  shear_refined=False,   # legacy alias of refined=1 (plate)
                  plot=True,             # writes <base>_mesh.png
                  boundary=None,         # "periodic" (default) | "aperiodic"
                  density=None,          # (n_mat,) beam mass rho
                  omega=None)            # user SG measure (3-D solid)
\end{lstlisting}

\noindent
and the argument that selects the physics is \texttt{n\_model}: it decides which
macro strain set $\Gamma_e$ is built. \texttt{n\_model}, \texttt{refined},
\texttt{boundary} and \texttt{omega} all default to \texttt{None}, meaning
\emph{take it from the YAML header}, so a self-describing SG file needs no
arguments at all --- that is what makes \texttt{opensg <yaml>} equivalent to
the call. The steps shared by all three macro models are:

\begin{enumerate}
\item \texttt{sg\_mesh.load\_sg\_input} parses the \texttt{.sc} or YAML into a
      dict carrying \texttt{dim} (the SG dimension, 1, 2 or 3), \texttt{nodes},
      \texttt{cells}, \texttt{mat\_id} and the materials. That \texttt{dim} is
      a \emph{parse result}, not something you write: the input contract has
      no \texttt{dim:} key and no \texttt{scale:} key, and both the dimension
      and the SG measure $\omega$ are read from the mesh itself. With
      \texttt{plot=True} it also writes \texttt{<base>\_mesh.png} of the actual
      mesh, elements coloured by material. Mixed element types in one SG are
      rejected --- the solver runs one homogeneous batch.
\item \texttt{sg\_mesh.\_cell\_basis} maps (SG dimension, nodes per element) to
      a basix cell type and basis degree; \texttt{fe\_jax.basis\_quadrature}
      then supplies the quadrature points \texttt{xi\_qp}, weights \texttt{W\_q},
      basis values \texttt{phi\_qn} and parametric gradients
      \texttt{dphi\_dxi\_qnp}. The quadrature degree is $6$ for higher-order
      elements and $2$ for linear ones.
\item \texttt{fe\_jax.setup.mesh\_to\_jax} produces the element-node coordinate
      array \texttt{x\_end} of shape $(E, N, d)$, and
      \texttt{mesh\_to\_periodic\_sparse\_assembly\_map} produces the master
      connectivity and the DOF map that carries periodicity. Periodicity rides
      entirely in the assembly map; no transformation matrix is formed. With
      \texttt{boundary="aperiodic"} the map is the identity and instead every
      bounding-box-face node gets $w=0$ Dirichlet; that option is available for
      the plate and solid models with \texttt{solver="direct"} only.
\item \texttt{sg\_materials.get\_heterogeneous\_C\_matrix} builds the per-element
      $6\times6$ table \texttt{C\_ess} of shape $(E,6,6)$, applying
      \texttt{angles} and \texttt{elem\_rotation} (per-element frames are
      honoured by all three models, not only the beam).
\end{enumerate}

\noindent
From there the routes diverge:

\begin{description}
\item[\texttt{n\_model=1}, beam.] Control goes to \texttt{\_beam\_homo\_kkt}:
      the four Euler--Bernoulli macro modes
      $[\varepsilon_{11}\;\kappa_1\;\kappa_2\;\kappa_3]$ are solved under four
      rigid-body Lagrange constraints, then the first-order chain runs and the
      generalized-Timoshenko reduction produces the $6\times6$. The result dict
      carries \texttt{C\_eff} (the Timoshenko $6\times6$) and
      \texttt{C\_eff\_EB} (the classical $4\times4$ from the same run). The
      output is \texttt{<base>.out}, named \emph{Timoshenko Beam} with
      \texttt{refined: 1} and \emph{Euler-Bernoulli Beam} with
      \texttt{refined: 0}, with $\omega$ in the banner. This path needs an SG
      of dimension $\ge2$.
\item[\texttt{n\_model=2}, plate.] The macro strain set is
      $\bar\varepsilon = [\varepsilon_{11}\;\varepsilon_{22}\;2\varepsilon_{12}\;
      \kappa_{11}\;\kappa_{22}\;\kappa_{12}]$ and the answer is the plate ABD,
      written as \emph{Classical Plate}. With \texttt{refined: 1} (or the
      legacy \texttt{shear\_refined=True}) the first-order warping ladder
      \texttt{plate\_shear\_ladder} runs as well, the result gains
      \texttt{G\_msg} ($2\times2$), \texttt{ABDG} (the $8\times8$ with the
      shear block on the diagonal) and \texttt{A6\_ladder}, and the file is
      named \emph{Reissner-Mindlin}.
\item[\texttt{n\_model=3}, solid.] The macro strain set is the six 3-D strains
      and the answer is the equivalent solid $6\times6$, named \emph{Cauchy
      Continuum} and written with the
      orthotropic-approximated engineering constants appended
      ($E_1$, $E_2$, $E_3$, $G_{12}$, $G_{13}$, $G_{23}$, $\nu_{12}$,
      $\nu_{13}$, $\nu_{23}$).
\end{description}

Models 2 and 3 share the assembly and solve code (\texttt{\_homo\_direct}), so
switching between a plate law and a solid law on the same mesh is a
one-character change. The shipped examples are numbered by capability under
\texttt{examples/OpenSG-solid/}: \texttt{1\_get\_plate\_props\_from\_1DSG}
(a stacking sequence in \texttt{layup\_db.yaml} $\rightarrow$
\texttt{layup\_db\_plate\_homo.out}, run with
\texttt{python 1d\_sg.py layup\_db.yaml}),
\texttt{3\_get\_plate\_props\_from\_2D\_SG},
\texttt{6\_get\_solid\_props\_from\_3D\_SG}, and
\texttt{7\_get\_beam\_props\_from\_SG}, whose script is:

\begin{lstlisting}[style=shellst]
cd examples/OpenSG-solid/7_get_beam_props_from_SG
python beam_homo_sg.py
\end{lstlisting}

\noindent
That one reads \texttt{RHC\_SW\_2UC\_45.yaml} (a two-unit-cell honeycomb
cross-section) with \texttt{n\_model = 1} and three materials overridden through
\texttt{material\_param} and \texttt{angles}, and writes
\texttt{RHC\_SW\_2UC\_45.out} plus \texttt{RHC\_SW\_2UC\_45\_mesh.png}.

\section{GPU readiness}
\label{sec:gpu}

The msg-solid engine has two interchangeable linear solvers, and the choice is
the \texttt{solver} argument of \texttt{plate\_homo\_2d}.

\texttt{solver="direct"} is the default. The jitted element kernels build the
sparse matrix, which is handed as a CSR to \texttt{\_sparse\_direct\_solve} ---
pypardiso (Intel MKL PARDISO) when it is importable, SciPy's SuperLU otherwise.
One factorization serves all macro strain columns. This is the fastest route on
a workstation, and it is the only stage of the calculation that leaves the
accelerator.

\texttt{solver="cg"} is the GPU route. The element kernels
(\texttt{calculate\_RHS\_and\_Ke\_batch\_periodic} and the jitted
\texttt{jacfwd} stiffness), the element-by-element operator
(\texttt{ebe\_jacobian\_product\_periodic}), the block and Chebyshev
preconditioners (\texttt{compute\_block\_inv\_diag},
\texttt{apply\_block\_precond}, \texttt{apply\_chebyshev\_precond},
\texttt{estimate\_max\_eigenvalue}) and the CG solver itself
(\texttt{fe\_jax.solve\_cg}) are pure JAX and run device-resident: no global
matrix is ever formed, so memory scales with the element batch rather than with
the bandwidth of an assembled operator. On a GPU host,
\texttt{plate\_homo\_2d(..., solver="cg")} therefore keeps the entire solve on
the device.

Switching solvers is digit-safe. The effective law
$\mathbf{C}_{\mathrm{eff}}$ is \emph{stationary} in $\mathbf{V}_0$ ---
equation~\eqref{eq:deff} evaluates the energy at its own minimiser --- so an
error $\delta$ in the fluctuation field enters the answer at
$O(\delta^2)$. A solver tolerance loose enough to be visible in $\mathbf{V}_0$
is still invisible in the printed matrix. Both routes are documented to produce
the same digits.

The msg-shell routes are one step behind. Their operators are assembled as
batched einsum contractions in the two-step form
$\mathbf{C}\mathbf{B}$ then $\mathbf{B}^{T}(\mathbf{C}\mathbf{B})$, and those
batches are written against the NumPy/\texttt{jnp} API and are portable to the
device; the sparse factorization (SuperLU for the ring and the 3-D shell SG,
pypardiso for the KKT solve in \texttt{msg\_solver}) is the remaining CPU stage.
For the SG sizes these routes handle --- a cross-section contour is a few
thousand DOF --- that factorization is not the bottleneck.
